\begin{abstract}
Laser ultrasonic testing enables non-contact inspection but suffers from low signal-to-noise ratio, requiring extensive signal averaging that reduces throughput and risks surface damage. Existing denoising methods inadequately preserve acoustic features critical for defect detection, limiting their practical utility. This paper proposes PMDF-Net, a physics-constrained multi-domain fusion network that jointly encodes time-domain waveforms and time-frequency representations to recover high-quality ultrasonic signals. The framework employs a dual-branch architecture that processes time-domain waveforms and time-frequency spectrograms in parallel, with a cross-domain fusion module that adaptively combines multi-scale features. The physics-constrained loss enforces consistency on arrival times, peak amplitudes, and waveform morphology through differentiable penalty functions, ensuring that acoustically relevant signal components are preserved. Training uses paired data consisting of low-averaging inputs (N=16 averages) and high-averaging targets (M=256 averages), enabling the network to learn the underlying signal characteristics from limited measurements. Experiments on laser ultrasonic benchmark datasets demonstrate that PMDF-Net achieves an SNR improvement of 11.4 dB while maintaining the best feature preservation quality among compared methods (CCC=0.981), and exceeds baseline F1-scores for defect detection. This enables a 16$\times$ reduction in averaging requirements while preserving equivalent defect detection performance compared to high-averaging references, substantially reducing inspection time and surface damage. Furthermore, zero-shot generalization is demonstrated from intact specimens to defect-containing workpieces without any fine-tuning. The proposed approach reduces laser ultrasonic inspection time and surface damage while preserving defect detection capability.
\end{abstract}